\documentclass[12pt]{article}
\usepackage{amsmath}
\usepackage{amsthm}
\usepackage{amssymb}
\usepackage{braket}
\usepackage{graphicx}

\title{Discrete Logarithm Algorithm in Diffie-Hellman Attacks:\\
A Quantum Computing Approach}
\author{}
\date{}

\begin{document}
\maketitle

\section{Introduction}
The Discrete Logarithm Algorithm (DLA) plays a crucial role in cryptanalysis, particularly in attacking the Diffie-Hellman key exchange protocol. This paper examines how DLA can be exploited to compromise the security of Diffie-Hellman and presents a quantum algorithm that provides a polynomial-time solution using the Quantum Fourier Transform (QFT).

\section{DLA in Diffie-Hellman Attack}
The Diffie-Hellman key exchange protocol relies on the computational hardness of the discrete logarithm problem. Given a prime modulus $p$, a generator $g$, and public values $g^a \bmod p$ and $g^b \bmod p$, an attacker who can efficiently solve the discrete logarithm problem can recover:

\begin{equation}
    a = \log_g(g^a \bmod p)
\end{equation}

With knowledge of $a$, the attacker can then compute the shared secret:

\begin{equation}
    k = (g^b)^a \bmod p
\end{equation}

The classical DLA has sub-exponential complexity, making it computationally infeasible for large prime moduli. However, quantum computing offers a polynomial-time solution.

\section{Quantum Algorithm for DLA}
We present a quantum algorithm based on the QFT that solves the discrete logarithm problem in polynomial time. The algorithm follows these steps:

\subsection{State Preparation}
Initialize two quantum registers in the state:

\begin{equation}
    \ket{\psi_0} = \frac{1}{\sqrt{p}} \sum_{x=0}^{p-1} \ket{x}\ket{0}
\end{equation}

\subsection{Function Evaluation}
Apply the modular exponentiation function:

\begin{equation}
    \ket{\psi_1} = \frac{1}{\sqrt{p}} \sum_{x=0}^{p-1} \ket{x}\ket{g^x \bmod p}
\end{equation}

\subsection{Quantum Fourier Transform}
Apply the QFT to the first register:

\begin{equation}
    QFT: \ket{x} \rightarrow \frac{1}{\sqrt{p}} \sum_{y=0}^{p-1} e^{2\pi ixy/p}\ket{y}
\end{equation}

The resulting state becomes:

\begin{equation}
    \ket{\psi_2} = \frac{1}{p} \sum_{x=0}^{p-1} \sum_{y=0}^{p-1} e^{2\pi ixy/p}\ket{y}\ket{g^x \bmod p}
\end{equation}

\subsection{Measurement}
Measure both registers. The period finding aspect of the QFT allows us to extract the discrete logarithm with high probability.

\section{Complexity Analysis}
The quantum algorithm achieves polynomial time complexity through:
\begin{enumerate}
    \item State preparation: $O(\log p)$ qubits and operations
    \item Modular exponentiation: $O(\log^3 p)$ operations
    \item QFT implementation: $O(\log^2 p)$ operations
\end{enumerate}

The total complexity is $O(\log^3 p)$, which is exponentially faster than the best classical algorithms.

\section{Conclusion}
The quantum solution to the discrete logarithm problem demonstrates the potential vulnerability of classical cryptographic systems to quantum attacks. This has significant implications for the security of Diffie-Hellman and similar protocols in a post-quantum world, necessitating the development of quantum-resistant alternatives.

\end{document}